\documentclass[11pt]{article}
\usepackage[T1]{fontenc}
\usepackage{lmodern,amsmath,amssymb,amsthm}
\usepackage[margin=1in]{geometry}
\usepackage[colorlinks=true,linkcolor=blue,citecolor=blue,urlcolor=blue]{hyperref}
\newtheorem{theorem}{Theorem}
\newtheorem{lemma}[theorem]{Lemma}
\newtheorem{proposition}[theorem]{Proposition}
\newcommand{\D}{\mathbb D}
\newcommand{\R}{\mathbb R}
\newcommand{\C}{\mathbb C}
\newcommand{\norm}[1]{\left\lVert#1\right\rVert}
\title{The leading constant for rational approximation of fractional inverse powers}
\author{Kaiwen Jin\\\small\href{mailto:yui.ui.siki@gmail.com}{yui.ui.siki@gmail.com}}
\date{13 October 2026}

\begin{document}
\maketitle
\begin{abstract}
For a fixed positive finite interval and a fixed exponent between zero and one,
we prove the Braess--Hackbusch conjecture for the leading constant in best
absolute uniform rational approximation of a fractional inverse power.
The approximants have numerator degree at most $n-1$ and denominator degree
at most $n$, with no restriction on poles outside the interval.
The proof combines strong asymptotics for Jacobi-weight multipoint Pad\'e
interpolants with a weighted finite-Blaschke construction on a slit disk.
Its alternating errors give the matching lower bound directly for the full
real rational class.
\end{abstract}

\section{The fixed-interval problem}

Rational approximations to fractional inverse powers are used in the
evaluation of fractional diffusion operators and related matrix functions.
An exponential convergence rate determines the growth of the required degree,
while its leading constant gives a more precise error prediction.
Braess and Hackbusch formulate the following constant as
Conjecture~4.8 in \cite{BH}.

For $0<a<b<\infty$ and $0<\gamma<1$, set
\begin{equation}\label{eq:error}
 E_{n-1,n}(a,b;\gamma)
 =\inf_{\substack{p,q\in\R[x],\ \deg p\le n-1,\ \deg q\le n\\
                  q(x)\ne0\ (x\in[a,b])}}
       \max_{x\in[a,b]}\left|x^{-\gamma}-\frac{p(x)}{q(x)}\right|.
\end{equation}
Let
\begin{equation}\label{eq:modulus}
 k=\frac{\sqrt b-\sqrt a}{\sqrt b+\sqrt a},\qquad
 K(k)=\int_0^{\pi/2}\frac{d\theta}{\sqrt{1-k^2\sin^2\theta}},\qquad
 \mu(k)=\frac{\pi K(\sqrt{1-k^2})}{2K(k)}.
\end{equation}
Thus $k$ is an elliptic modulus, rather than its square.

\begin{theorem}\label{thm:main}
For every fixed $0<a<b<\infty$ and every fixed $0<\gamma<1$,
\begin{equation}\label{eq:main}
 \lim_{n\to\infty}e^{2n\mu(k)} E_{n-1,n}(a,b;\gamma)
 =4(ab)^{-\gamma/2}\sin(\pi\gamma).
\end{equation}
Here $E_{n-1,n}$ is the absolute uniform error in \eqref{eq:error}, and
$k$ and $\mu$ are defined in \eqref{eq:modulus}.
\end{theorem}

The interval and exponent remain fixed in this limit. In particular,
\eqref{eq:main} is not a uniform statement at $\gamma=0,1$ or for intervals
approaching the branch point or becoming unbounded.

The two contributions to the constant have different origins.
The jump of the fractional power, evaluated through the Jacobi-weight
interpolation theorem of Baratchart and Yattselev \cite{BY}, gives the factor
$2\sin(\pi\gamma)x^{-\gamma}$ in the interpolation error.
A weighted Blaschke product then makes the remaining error nearly constant
in magnitude on the interval. Its geometric mean supplies
$(ab)^{-\gamma/2}$, and its two boundary values supply the second factor~2.
We prove the needed weighted interval construction below. The annulus
geometry used in this construction is the classical geometry recalled in
\cite[Section~3]{BPS}; none of its meromorphic approximation conclusions is
used. The elliptic normalization is fixed by the relative square-root
benchmark \cite[Theorem~4.2]{BH}.

\section{Weighted Blaschke products on an interval}

Write $\D=\{z:|z|<1\}$ and fix $I=[u,v]$ with $0<u<v<1$.
A real monic finite Blaschke product with $m$ zeros in $I$ has the form
\begin{equation}\label{eq:blaschke}
 B_m(z)=\prod_{j=1}^m\frac{z-e_j}{1-e_jz},\qquad e_j\in I.
\end{equation}
For a positive continuous weight $W$ on $I$, define
\[
 T_m(W;I)=\inf_{e_1,\ldots,e_m\in I}\max_{z\in I}W(z)|B_m(z)|.
\]
The weight will later be a fractional inverse power in a half-plane
coordinate. It must be retained to determine the leading constant.

The slit disk $\D\setminus I$ is conformally equivalent to an annulus
$\rho<|\zeta|<1$, with $0<\rho<1$. Choose the real-symmetric map
$\Phi$ so that its boundary values at $v$ and $u$ are $\rho$ and $-\rho$,
respectively. This is the restriction of the normalized doubled-domain
map in \cite[Section~3]{BPS}. Let $Z=\Phi^{-1}$, including its boundary
values on the two sides of the slit, and let $\lambda_I$ be the pushforward
of $d\theta/(2\pi)$ under $Z(\rho e^{i\theta})$.

We recall two geometric facts used in the proof. Reflection across the slit
continues the inverse map by
\begin{equation}\label{eq:inverse-reflection}
 Z(\rho^2/\zeta)=Z(\zeta).
\end{equation}
The inverse is analytic in the inner-circle parameter, with a quadratic
critical point at each endpoint angle. To see the endpoint statement,
open the slit angle $2\pi$ by the coordinate $\eta=\sqrt{z-v}$, or
$\sqrt{z-u}$ at the other endpoint. Schwarz reflection makes the map
analytic in $\eta$. A zero of order at least two in this coordinate would
map a half-neighborhood to an angle at least $2\pi$, contrary to the
image half-neighborhood. Its derivative is therefore nonzero, and the
inverse has quadratic behavior in the original coordinate. Away from the
endpoints, ordinary reflection across analytic arcs applies
\cite[Section~7.2]{BN}.

The measure $\lambda_I$ is the normalized disk-Green equilibrium measure.
Indeed $-\log|\Phi|$ is zero on the unit circle and equals $-\log\rho$ on
the slit, so it equals the equilibrium potential stated in
\cite[Section~3]{BPS}, by the maximum principle.
The inward normal flux along both sides of the slit is the pushforward
of $d\theta$; division by $2\pi$ gives mass one. The quadratic endpoint
behavior makes the normal derivative integrable, and the fluxes on small
endpoint circles tend to zero. This identifies the angular pushforward
with the equilibrium measure. Equivalently, one can apply Green's identity
to the disk kernel $\log|(1-tz)/(z-t)|$, whose point-source flux is $2\pi$.

\begin{samepage}
\begin{lemma}\label{lem:weighted}
For every fixed positive continuous $W$ on $I$,
\begin{equation}\label{eq:weighted-minimum}
 T_m(W;I)=2\rho^m C_W(1+o(1)),\qquad
 C_W=\exp\left(\int_I\log W\,d\lambda_I\right).
\end{equation}
Moreover, for every fixed $\varepsilon>0$ there is a sequence of products
\eqref{eq:blaschke} with distinct interior zeros whose normalized zero counting
measures $m^{-1}\sum_{j=1}^m\delta_{e_j}$ converge to $\lambda_I$,
with weighted norms at most
$2C_W e^{2\varepsilon}\rho^m(1+o(1))$ and with $m+1$ alternating weighted
values of magnitudes at least
$2C_W e^{-2\varepsilon}\rho^m(1+o(1))$.
\end{lemma}
\end{samepage}

\begin{proof}
We first take $W=e^P$, where $P$ is a real polynomial. Put
\[
 C=\exp\left(\frac1{2\pi}\int_0^{2\pi}
 P(Z(\rho e^{i\theta}))\,d\theta\right),\qquad
 g(\theta)=\log C-P(Z(\rho e^{i\theta})).
\]
The function $g$ is real analytic, even and of mean zero.
Its cosine expansion $g(\theta)=\sum_{j\ge1}g_j\cos(j\theta)$ has
exponentially decaying real coefficients. Define
\begin{equation}\label{eq:compensator}
 F(\zeta)=\sum_{j\ge1}\frac{g_j}{\rho^j-\rho^{-j}}
                  (\zeta^j-\zeta^{-j}),\qquad A(\zeta)=e^{F(\zeta)}.
\end{equation}
The series converges normally in a neighborhood of the closed annulus.
Its real parts give
\begin{equation}\label{eq:factor-moduli}
 |A(\zeta)|=1\quad(|\zeta|=1),\qquad
 |A(\rho e^{i\theta})|=C e^{-P(Z(\rho e^{i\theta}))}.
\end{equation}
There is a single-valued logarithm: the vanishing mean has removed a
possible logarithmic term on the annulus.

Continue $F$ to the reflected annulus using
\begin{equation}\label{eq:factor-reflection}
 F(\rho^2/\zeta)=2\{\log C-P(Z(\zeta))\}-F(\zeta),
                  \qquad \rho<|\zeta|<1.
\end{equation}
On the inner circle this formula is the complex conjugate of $F$, so it
agrees with its local analytic continuation. The continued $A$ is
nonzero and bounded on the closed doubled annulus. Here using an
exponential polynomial weight is useful: $P$ is analytic and bounded on
the entire closed disk, including the outer boundary.

The function
\begin{equation}\label{eq:symmetrization}
 H_m(\zeta)=\zeta^m A(\zeta)
       +\rho^{2m}\zeta^{-m}A(\rho^2/\zeta)
\end{equation}
is invariant under $\zeta\mapsto\rho^2/\zeta$.
It therefore descends to a holomorphic function $h_m$ in $\D$.
At a slit endpoint the sheet exchange changes the sign of the local
square-root coordinate, so the invariant expression has only even terms
and its endpoint singularity is removable. It is also continuous, and
analytic locally, across the outer circle.

On that circle the first term in \eqref{eq:symmetrization} has modulus one,
and the second is bounded by $M\rho^{2m}$, with $M$ independent of $m$.
Consequently
\begin{equation}\label{eq:outer-bound}
 \big||h_m|-1\big|\le M\rho^{2m}\quad\hbox{on }\partial\D.
\end{equation}
For large $m$ the boundary winding is $m$: $\Phi$ has boundary degree one,
$A=e^F$ has winding zero, and the second term is a smaller perturbation.
The argument principle gives exactly $m$ zeros in the disk, counted with
multiplicity.

Set $\psi(\theta)=\operatorname{Im}F(\rho e^{i\theta})$.
On the slit, \eqref{eq:symmetrization} becomes
\begin{equation}\label{eq:cosine}
 e^{P(Z(\rho e^{i\theta}))}h_m(Z(\rho e^{i\theta}))
       =2C\rho^m\cos(m\theta+\psi(\theta)).
\end{equation}
On the upper-side parametrization $0\le\theta\le\pi$, the real Laurent
coefficients imply $\psi(0)=\psi(\pi)=0$. Its derivative is bounded.
Thus $m\theta+\psi(\theta)$ is strictly increasing for large $m$, from
zero to $m\pi$. Equation \eqref{eq:cosine} gives $m$ distinct interior
zeros and $m+1$ alternating values $\pm2C\rho^m$, including the endpoints.
These are all the disk zeros by the preceding count.

Factor $h_m=B_m O_m$, where $B_m$ is the real monic product of these zeros.
The factor $O_m$ is holomorphic and zero-free in the disk. Since $|B_m|=1$
on its boundary, the maximum principle applied to $O_m$ and $1/O_m$ gives
\begin{equation}\label{eq:outer-factor}
 1-M\rho^{2m}\le|O_m(z)|\le1+M\rho^{2m}\quad(z\in\overline\D).
\end{equation}
It is positive on $I$: both factors are real there and
$h_m(v)>0$, $B_m(v)>0$. Hence the weighted norm of $B_m$ is at most
$2C\rho^m/(1-M\rho^{2m})$, while its alternating values have magnitudes
at least $2C\rho^m/(1+M\rho^{2m})$.

These values give the matching lower bound for all other products.
If another product had smaller norm than all their magnitudes, its
difference from $B_m$ would have at least $m$ distinct roots in $I$.
Write the products as $p/p^*$ and $q/q^*$, with
$p^*(z)=z^m p(1/z)$ and $q^*(z)=z^m q(1/z)$.
The difference numerator $N=pq^*-qp^*$ has degree at most $2m$ and obeys
\[
 z^{2m}N(1/z)=-N(z).
\]
The $m$ roots in $I\subset(0,1)$ therefore have $m$ distinct reciprocal
roots outside the disk. Also $N(1)=0$. This gives $2m+1$ distinct roots,
which is impossible for a nonzero numerator. Identical products cannot
have a strictly smaller norm. This proves \eqref{eq:weighted-minimum}
for $W=e^P$.

The zero angles satisfy
$m\theta_{j,m}+\psi(\theta_{j,m})=(j+\tfrac12)\pi$ for $0\le j<m$.
They differ from the equally spaced midpoint angles by $O(1/m)$.
Their normalized zero counting measures consequently converge to $\lambda_I$.

For a general positive continuous weight choose $P$ with
$\norm{\log W-P}_{\infty,I}<\varepsilon$.
This approximation can be obtained from Bernstein polynomials: on
$[0,1]$ their binomial weights have centered second moment
$t(1-t)/N$, so for continuous $f$ the approximation error is at most
$\omega_f(\delta)+2\norm f_\infty/(4N\delta^2)$.
Taking $\delta=N^{-1/4}$ proves uniform convergence.
The norms for $W$ and $e^P$ differ by at most $e^{\pm\varepsilon}$,
as do their geometric means. The constructed products therefore have
the bounds in the lemma. The same comparison sandwiches the minimum.
First let $m$ tend to infinity with $P$ fixed, and then let
$\varepsilon$ decrease to zero, proving \eqref{eq:weighted-minimum}.
\end{proof}

We now express the lemma in the half-plane coordinates used for fractional
powers. Let $J=[\alpha,\beta]$, $0<\alpha<\beta$, and choose $0<d<\alpha$.
The disk coordinate $u(s)=(s-d)/(s+d)$ gives the exact identity
\begin{equation}\label{eq:half-plane-product}
 \frac{u(s)-u(t)}{1-u(s)u(t)}=\frac{s-t}{s+t}.
\end{equation}
The half-plane Green kernel on positive real points is
$\log|(s+t)/(s-t)|$. Simultaneous inversion by $\alpha\beta$ leaves this
kernel and $J$ invariant. Uniqueness of the normalized equilibrium
measure $\lambda_J$ thus gives
\begin{equation}\label{eq:log-mean}
 \int_J\log s\,d\lambda_J(s)=\frac12\log(\alpha\beta).
\end{equation}
For the weight $s^{-2\gamma}$ its geometric mean is consequently
$(\alpha\beta)^{-\gamma}$.

\begin{lemma}\label{lem:elliptic}
For $J=[\sqrt a,\sqrt b]$, the annulus radius in
Lemma~\ref{lem:weighted} is $\rho=e^{-\mu(k)}$, with the convention
\eqref{eq:modulus}: $k=(\sqrt b-\sqrt a)/(\sqrt b+\sqrt a)$ and
$\mu(k)=\pi K(\sqrt{1-k^2})/(2K(k))$.
\end{lemma}
\begin{proof}
Take the unweighted products constructed in Lemma~\ref{lem:weighted}
and put $m=2n$. In the half-plane coordinate they have norm
$2\rho^{2n}(1+o(1))$ and $2n+1$ alternating values with this leading
magnitude. Write $B(s)=P(s)/P(-s)$, where $P$ is a degree-$2n$
polynomial whose roots lie in $J$, and set
\[
 v_n(x)=\sqrt x\frac{1+B(\sqrt x)}{1-B(\sqrt x)}.
\]
The even numerator and odd denominator in this expression show that
$v_n$ has type $(n,n-1)$ in $x$. Its denominator is nonzero on $[a,b]$
for large $n$, since $\norm B_{\infty,J}<1$.
Its signed relative error is $(\sqrt x-v_n(x))/\sqrt x=-2B/(1-B)$,
with norm at most
$4\rho^{2n}(1+o(1))$ and with $2n+1$ alternating values of at least that
leading magnitude.

Any type $(n,n-1)$ rational function with a smaller relative error than
all those values would differ from $v_n$ at $2n$ distinct zeros in the
interval. Its difference numerator has degree at most $2n-1$ and its
denominator is nonzero there, a contradiction. Hence the best relative
square-root error is $4\rho^{2n}(1+o(1))$.
Theorem~4.2 of \cite{BH} gives the same error as
$4e^{-2n\mu(k)}(1+o(1))$. Taking $2n$-th roots proves the claimed identity.
\end{proof}

\section{Jacobi interpolation and the absolute error}

Fix $c=a/2$ and use
\begin{equation}\label{eq:mobius}
 z=\frac{x+c}{x-c},\qquad
 M_\gamma(z)=\frac{(z-1)^{\gamma-1}}{(z+1)^\gamma},\qquad
 x^{-\gamma}=c^{-\gamma}(z-1)M_\gamma(z).
\end{equation}
The transformed interval is fixed and strictly to the right of~1.
The branch of $M_\gamma$ off $[-1,1]$ is determined by its value $1/z$
at infinity. It has the Markov representation
\begin{equation}\label{eq:jacobi}
 M_\gamma(z)=\int_{-1}^1\frac{d\nu_\gamma(t)}{z-t},\qquad
 d\nu_\gamma(t)=\frac{\sin(\pi\gamma)}{\pi}
              (1-t)^{\gamma-1}(1+t)^{-\gamma}\,dt.
\end{equation}
For an explicit verification, set $v=(1-t)/(1+t)$. The integral of the
weight divided by $z-t$ becomes
\[
 \int_0^\infty\frac{v^{\gamma-1}}{(z-1)+(z+1)v}\,dv.
\]
Scaling $v$ by $(z-1)/(z+1)$ and using the beta integral gives
\eqref{eq:jacobi}. That beta integral follows also from the second
fractional-power identity in \cite[equation~(3.5)]{BH} at $x=1$.
Both endpoints are integrable because $0<\gamma<1$.

\subsection{The interpolation formula}

We state the particular consequence of \cite[Theorem~2]{BY} needed here.
For a triangular scheme of $2n$ distinct nodes $\xi_{j,n}\in[a,b]$ whose
transformed empirical measures converge weakly, let $r_n$ be the
type $(n-1,n)$ interpolant to $x^{-\gamma}$ at those nodes, obtained from
the Markov interpolant to $M_\gamma$ by \eqref{eq:mobius}.
Then
\begin{equation}\label{eq:interpolation-error}
 x^{-\gamma}-r_n(x)
 =\{2\sin(\pi\gamma)+o(1)\}x^{-\gamma}
      \prod_{j=1}^{2n}\frac{\sqrt x-\sqrt{\xi_{j,n}}}
                           {\sqrt x+\sqrt{\xi_{j,n}}},
\end{equation}
uniformly on $[a,b]$. We verify its type, admissibility and constant,
as well as the source assumptions.

Let $z_j=(\xi_{j,n}+c)/(\xi_{j,n}-c)>1$ and
$V(z)=\prod_{j=1}^{2n}(z-z_j)$. The measure $d\nu_\gamma/V$ on $[-1,1]$
is positive and integrable. Its positive definite Gram matrix yields
a real monic orthogonal polynomial $q_n$ of degree $n$.
All its roots are simple in $(-1,1)$: if it had fewer than $n$
sign-changing roots, multiplying by the polynomial of just those roots
would give a nonzero constant-sign integrand against the positive
measure, contradicting orthogonality.
Define
\[
 R_n(z)=\int_{-1}^1\frac{q_n(t)}{V(t)(z-t)}\,d\nu_\gamma(t),\qquad
 p_n(z)=q_n(z)M_\gamma(z)-V(z)R_n(z).
\]
Difference quotients show that $p_n$ is a polynomial. For powers $z^j$
with $j\ge n$, the coefficient of $(V(z)-V(t))/(z-t)$ has degree at most
$2n-1-j\le n-1$ in $t$, and its integral against
$q_n(t)d\nu_\gamma(t)/V(t)$ is zero. Hence $\deg p_n\le n-1$.
The quotient $\Pi_n=p_n/q_n$ interpolates $M_\gamma$ at the $z_j$,
vanishes at infinity, and is the Pad\'e approximant used in \cite{BY}:
orthogonality gives $R_n(z)=O(z^{-n-1})$, so that
$(q_nM_\gamma-p_n)/V=O(z^{-n-1})$ as required there.
Its uniqueness follows by counting $2n$ interpolation zeros of a
difference numerator of degree at most $2n-1$.

Multiplying $\Pi_n(z(x))$ by $c^{-\gamma}(z(x)-1)$ gives exactly type
$(n-1,n)$ in $x$: after multiplying its numerator and denominator by
$(x-c)^n$, the single leftover numerator factor $x-c$ cancels
$z-1=2c/(x-c)$.
The roots of $q_n$ give strictly negative original poles
$x=c(z+1)/(z-1)$. The resulting $r_n$ is therefore admissible for
\eqref{eq:error}.

To evaluate the source formula, write
\[
 \mathfrak w(z)=\sqrt{z^2-1}\sim z,\qquad
 \eta(z)=\frac1{z+\mathfrak w(z)},\qquad
 d\omega(t)=\frac{dt}{\pi\sqrt{1-t^2}}.
\]
The density of $\nu_\gamma$ relative to $\omega$ is
\[
 \dot\nu_\gamma(t)=\sin(\pi\gamma)
              \left(\frac{1-t}{1+t}\right)^{\gamma-1/2}.
\]
Its geometric mean is $G=\sin(\pi\gamma)$, since the arcsine means of
$\log(1-t)$ and $\log(1+t)$ agree. Its normalized Szeg\H{o} function is
\begin{equation}\label{eq:szego}
 S(z)=\left(\frac{1-\eta(z)}{1+\eta(z)}\right)^{\gamma-1/2}.
\end{equation}
This normalization can be checked directly in \cite[equation~(2.11)]{BY}:
for $t=\cos\theta$, the cosine coefficient of
$\log((1-t)/(1+t))$ is $-4/j$ for odd $j$ and zero for even $j$.
The identity
\[
 \frac{\mathfrak w(z)}{z-\cos\theta}
       =1+2\sum_{j\ge1}\eta(z)^j\cos(j\theta)
\]
then gives $\log S=-2(\gamma-\tfrac12)
\sum_{j\text{ odd}}\eta^j/j$, which is \eqref{eq:szego}.
The endpoint logarithms are integrable, and the series is normally
convergent on compact subsets off the cut, justifying the integration.

Theorem~2 of \cite{BY}, in the notation of its arXiv version
0911.3850v1, now states
\begin{equation}\label{eq:by-formula}
 (M_\gamma-\Pi_n)\mathfrak w
           =[2G+o(1)]S^2 b_n,\qquad
 b_n(z)=\prod_{j=1}^{2n}
          \frac{\eta(z)-\eta(z_j)}{1-\eta(z)\eta(z_j)}.
\end{equation}
The Jacobi exponents are $\gamma-1,-\gamma\in(-1,0)$.
The remaining density factor is positive and constant, so its logarithm
is in every smoothness class in the source. Taking its smoothness
parameter $s=3$ satisfies \cite[equation~(2.9)]{BY}, namely
\[
 \gamma-1,-\gamma\in(-s,s)\cap(-1,\infty),\qquad s=3.
\]
All nodes remain in a fixed compact interval to the right of~1.
On the cut each factor has modulus one, and on any compact subset off
the cut each disk Blaschke factor has modulus bounded strictly below
one, uniformly in the nodes. The product therefore decays geometrically.
Together with weak convergence of the node measures, these are exactly
the three conditions of Definition~1 in \cite{BY}.

Finally,
\[
 \eta(z(x))=\frac{\sqrt x-\sqrt c}{\sqrt x+\sqrt c},\qquad
 \left(\frac{1-\eta}{1+\eta}\right)^2
       =\frac{z-1}{z+1}=\frac c x.
\]
Equation \eqref{eq:half-plane-product} turns $b_n$ into the product in
\eqref{eq:interpolation-error}. Also
\[
 \frac{c^{-\gamma}(z-1)}{\mathfrak w(z)}S(z)^2
       =c^{-\gamma}\left(\frac{z-1}{z+1}\right)^{\gamma}
       =x^{-\gamma}.
\]
Multiplying \eqref{eq:by-formula} by the remaining factor proves
\eqref{eq:interpolation-error}, with a uniform remainder on the fixed
positive interval.

\subsection{Completion of the minimax bounds}

\begin{proof}[Proof of Theorem~\ref{thm:main}]
Fix $\varepsilon>0$. Apply Lemma~\ref{lem:weighted} on
$J=[\sqrt a,\sqrt b]$ in the disk coordinate
$u(s)=(s-\sqrt c)/(s+\sqrt c)$, with weight $s^{-2\gamma}$ and
degree $m=2n$.
By \eqref{eq:log-mean} its geometric mean is
$C=(ab)^{-\gamma/2}$.
The resulting products have weighted norms at most
$2Ce^{2\varepsilon}\rho^{2n}(1+o(1))$ and $2n+1$ alternating weighted
values with magnitudes at least
$2Ce^{-2\varepsilon}\rho^{2n}(1+o(1))$.

Their zeros $t_{j,n}$ yield interpolation nodes $\xi_{j,n}=t_{j,n}^2$.
The normalized counting measures of these nodes, including after the fixed
M\"obius transformation, converge because the zero angles are midpoint
angles plus $O(1/n)$. Thus \eqref{eq:interpolation-error} applies to this
particular scheme. It gives an admissible interpolant with norm at most
\[
 4C\sin(\pi\gamma)e^{2\varepsilon}\rho^{2n}(1+o(1)).
\]
The same error has $2n+1$ alternating values of magnitudes at least
\[
 4C\sin(\pi\gamma)e^{-2\varepsilon}\rho^{2n}(1+o(1)),
\]
since its multiplier in \eqref{eq:interpolation-error} is positive at
these points for large $n$.

These values bound the error of every rational function in
\eqref{eq:error}. Otherwise its difference from the interpolant would
have $2n$ distinct zeros in $[a,b]$. The difference numerator has degree
at most $2n-1$, and both denominators are nonzero on the interval, a
contradiction. This argument includes lower actual degrees and places
no condition on poles outside $[a,b]$.

Using $\rho=e^{-\mu(k)}$ from Lemma~\ref{lem:elliptic}, we obtain
\[
 4C\sin(\pi\gamma)e^{-2\varepsilon}
 \le\liminf_n e^{2n\mu(k)}E_{n-1,n}
 \le\limsup_n e^{2n\mu(k)}E_{n-1,n}
 \le4C\sin(\pi\gamma)e^{2\varepsilon}.
\]
All choices depend only on the fixed interval, exponent and
$\varepsilon$. Letting $\varepsilon$ tend to zero proves
\eqref{eq:main}.
\end{proof}

\clearpage
\begingroup
\small
\begin{thebibliography}{9}
\setlength{\itemsep}{2pt}
\bibitem{BN}
J. Bak and D. J. Newman,
\emph{Complex Analysis}, third edition, Springer, 2010.
\href{https://doi.org/10.1007/978-1-4419-7288-0}{doi:10.1007/978-1-4419-7288-0}.

\bibitem{BPS}
L. Baratchart, V. A. Prokhorov and E. B. Saff,
Asymptotics for minimal Blaschke products and best $L_1$ meromorphic
approximants of Markov functions,
\emph{Computational Methods and Function Theory} \textbf{1} (2001), 501--520.
\href{https://math.vanderbilt.edu/saffeb/texts/193.pdf}{Author preprint}.

\bibitem{BY}
L. Baratchart and M. Yattselev,
Convergent interpolation to Cauchy integrals over analytic arcs with
Jacobi-type weights,
\emph{International Mathematics Research Notices} \textbf{2010} (2010),
no.~22, 4211--4275.
\href{https://doi.org/10.1093/imrn/rnq026}{doi:10.1093/imrn/rnq026}.
\href{https://arxiv.org/abs/0911.3850v1}{Preprint arXiv:0911.3850v1}.

\bibitem{BH}
D. Braess and W. Hackbusch,
On the approximation of Stieltjes functions by exponential sums and rational
functions with applications to partial differential equations,
\emph{Numerische Mathematik} \textbf{158} (2026), 455--490.
\href{https://doi.org/10.1007/s00211-025-01523-1}{doi:10.1007/s00211-025-01523-1}.
\end{thebibliography}
\endgroup

\paragraph*{AI-use statement.}
OpenAI Codex assisted with literature retrieval, mathematical derivations,
proof checking and preparation of the manuscript.
Independent mathematical review is requested with the submission.
\end{document}
